\documentclass[dvipsnames]{article}

\usepackage[utf8]{inputenc}
\usepackage{fontspec}
\usepackage{polyglossia}
\usepackage{unicode-math}
\usepackage{csquotes}
\setdefaultlanguage{russian}
\setotherlanguage{english}

\setmainfont{CMU Serif}
\setsansfont{CMU Sans Serif}
\setmonofont{CMU Typewriter Text}
% \setmathfont{Latin Modern Math}

\usepackage[letterpaper,top=2cm,bottom=2cm,left=3cm,right=3cm,marginparwidth=1.75cm]{geometry}

\usepackage{amsmath}
\usepackage{graphicx}
\usepackage{amsthm}
\usepackage{array}
\usepackage{fancyhdr}
\usepackage{bussproofs}
\usepackage[colorlinks=true, allcolors=blue]{hyperref}
\usepackage[backend=biber]{biblatex}
\usepackage{xcolor}

\usepackage{minted2}
\usemintedstyle{xcode}
\setminted{escapeinside=||}
\newminted[stella]{cpp}{escapeinside=||,frame=single}
\newmintinline[haskell]{haskell}{fontsize=\scriptsize}
\newmintinline[java]{java}{fontsize=\scriptsize}
\newmintinline[python]{python}{fontsize=\scriptsize}
\newmintinline[code]{c}{fontsize=\small}
\newmintinline[smallcode]{c}{fontsize=\small}
\let\smallcodeOrig\smallcode
\renewcommand{\smallcode}[2][]{%
  \ifmmode
    \if\relax\detokenize{#1}\relax
      \mathchoice{\mbox{\normalfont\small\ttfamily #2}}{\mbox{\normalfont\small\ttfamily #2}}{\mbox{\normalfont\small\ttfamily #2}}{\mbox{\normalfont\small\ttfamily #2}}%
    \else
      \mathchoice{\mbox{\normalfont\small\ttfamily #2}}{\mbox{\normalfont\small\ttfamily #2}}{\mbox{\normalfont\small\ttfamily #2}}{\mbox{\normalfont\small\ttfamily #2}}%
    \fi
  \else
    \smallcodeOrig[#1]{#2}%
  \fi}
\let\codeOrig\code
\renewcommand{\code}[2][]{%
  \ifmmode
    \if\relax\detokenize{#1}\relax
      \mathchoice{\mbox{\normalfont\small\ttfamily #2}}{\mbox{\normalfont\small\ttfamily #2}}{\mbox{\normalfont\small\ttfamily #2}}{\mbox{\normalfont\small\ttfamily #2}}%
    \else
      \mathchoice{\mbox{\normalfont\small\ttfamily #2}}{\mbox{\normalfont\small\ttfamily #2}}{\mbox{\normalfont\small\ttfamily #2}}{\mbox{\normalfont\small\ttfamily #2}}%
    \fi
  \else
    \codeOrig[#1]{#2}%
  \fi}

\usepackage{sourcecodepro}
\usepackage{mdframed}

\DeclareMathOperator{\xor}{\mathsf{xor}}
\DeclareMathOperator{\tru}{\mathsf{tru}}
\DeclareMathOperator{\fls}{\mathsf{fls}}

\addbibresource{../references.bib}

\begin{document}

\pagestyle{fancy}
\fancyhf{}
\chead{Задание выполнил \textcolor{blue}{Бенжамин Пирс (замените на своё имя)}}
\cfoot{Языки программирования. Семантика и системы типов (весна 2026, Университет ИТМО)}

\section*{Тема 8. Теоретическое задание \textcolor{blue}{(выполнил Бенжамин Пирс)}}%

\begin{enumerate}
  \item Для каждого из следующих типов выпишите \textbf{все} подтипы\footnote{В случае программной генерации подтипов, программа должна быть написана вами и прикреплена к решению.}:
    \begin{enumerate}
      \item $\code{\{ a : <b : Nat> → Nat \} → <e : <f : Bool> >}$
      \item $\code{List[\{a : Nat, b : Bool\}] → List[<a : Nat, b : Bool>]}$
      \item $\code{(<a : Ref \{b : Top\} → Bot> → Top) → Bot}$
    \end{enumerate}

  \color{blue}
  \textbf{Решение.}
  \emph{место для решения}
  \color{black}

  \clearpage

  \item Докажите \cite[Лемма~16.1.2]{TaPL-ru} для типизированного $\lambda$-исчисления с подтипами~\cite[\S 15]{TaPL-ru}, записями
    (используя правило \textsc{S-Rcd} вместо трёх отдельных правил),
    списками, ссылками и типом \code{Top} (без \code{Bot}):
    \begin{enumerate}
      \item $\code{S <: S}$ можно вывести без использования \textsc{S-Refl};
      \item если $\code{S <: T}$ выводимо, то оно выводимо без использования \textsc{S-Trans}.
    \end{enumerate}

  \color{blue}
  \textbf{Решение.}
  \emph{место для решения}
  \color{black}

  \clearpage

  \item Покажите \cite[Лемма~16.3.2]{TaPL-ru}, что для любых двух типов $\code{T$_1$}$ и $\code{T$_2$}$
    в типизированном $\lambda$-исчислении с подтипами~\cite[\S 15]{TaPL-ru}, записями (правило \textsc{S-Rcd}~\cite[\S 16.1,~стр.~227]{TaPL-ru} вместо трёх отдельных правил) и типом \code{Top}
    (без \code{Bot}):
    \begin{enumerate}
      \item возможно построение объединения $\code{J(T$_1$, T$_2$)}$, такого что
      $\code{T$_1$ <: J}$, $\code{T$_2$ <: J}$ и для всех типов $\code{U}$, если $\code{T$_1$ <: U}$ и $\code{T$_2$ <: U}$,
      то $\code{J <: U}$;
      \item возможно построение пересечения $\code{M(T$_1$, T$_2$)}$, такого что
      $\code{M <: T$_1$}$, $\code{M <: T$_2$}$ и для всех типов $\code{U}$, если $\code{U <: T$_1$}$ и $\code{U <: T$_2$}$,
      то $\code{U <: M}$.
    \end{enumerate}

  \color{blue}
  \textbf{Решение.}
  \emph{место для решения}
  \color{black}
\end{enumerate}

\section*{Список литературы}
\printbibliography[heading=none]

\end{document}

